%==============================================================================
%  Invariance vs robustness -- INTERNAL reader script (not for sharing)
%  Companion to pitch_deck.pdf. One block per slide: the exact words to say,
%  plus a backup bank of anticipated questions with short answers.
%==============================================================================
\documentclass[11pt]{article}
\usepackage[T1]{fontenc}
\usepackage[utf8]{inputenc}
\usepackage{lmodern}
\usepackage[a4paper,margin=2.3cm]{geometry}
\usepackage{amsmath,amssymb}
\usepackage{xcolor}
\usepackage{enumitem}
\usepackage{titlesec}
\usepackage{parskip}

\definecolor{TUMblue}{HTML}{0065BD}
\definecolor{TUMdarker}{HTML}{003359}
\definecolor{slate}{HTML}{52606D}

\newcommand{\etaL}{\ensuremath{\eta/L}}

\titleformat{\section}{\large\bfseries\color{TUMdarker}}{\thesection.}{0.5em}{}
\titlespacing{\section}{0pt}{1.4em}{0.5em}

% one slide block: \slide{n}{title}{script}
\newcommand{\slide}[3]{%
  \par\vspace{0.9em}%
  {\color{TUMblue}\bfseries Slide #1 \textbullet\ #2}\par
  \vspace{0.25em}#3\par}

\newcommand{\qa}[2]{\par\vspace{0.5em}{\bfseries #1}\ \ #2\par}

\begin{document}

\begin{center}
  {\LARGE\bfseries\color{TUMdarker} Invariance vs robustness}\\[0.3em]
  {\large\color{slate} Internal reader script and question bank}\\[0.4em]
  {\small\color{slate} Companion to \texttt{pitch\_deck.pdf}. Not for sharing.
  Nineteen slides; script per slide, then the backup question bank.}
\end{center}
\vspace{0.6em}
\hrule
\vspace{0.4em}

{\small\color{slate} How to use this: each block below matches one slide in
\texttt{pitch\_deck.pdf}, in order. The text in each block is what to say out loud on that
slide. It is written for a pitch to one person (Maha), not a conference talk, so it is short and
plain. The question bank at the end is not in the shared deck; keep it here.}

\section{Per-slide script}

\slide{1}{Title --- ``Invariance vs robustness''}{%
This is a project that grew out of the practical course. It is about two properties of a vision
model, invariance and adversarial robustness. The literature has related them and reached opposite
conclusions. I will show what actually orders robustness across models, and it is not the invariance
metric people track.}

\slide{2}{The two properties, on one image}{%
Here is one image and a standard ImageNet classifier. On top, a change of one part in 255, which you
cannot see, turns a heron into an elephant. That is adversarial robustness, or the lack of it. On the
bottom, sliding the image around does nothing to the label. That is shift invariance. The whole
question is how these two relate.}

\slide{3}{The literature relates them, and disagrees}{%
One line of work, anti-aliasing and equivariance, reports that invariance helps robustness: Zhang,
Saha and Gokhale, Grabinski, Wang. Another, Ge and Kamath, reports the opposite, or a proven
trade-off. They study different data, models, and attacks. And, importantly, the ``helps'' side only
ever used gradient attacks like FGSM and PGD, never a strong parameter-free attack. So the
disagreement is not actually resolved. My question is: is there a quantity that orders robustness
across models?}

\slide{4}{How we frame the project}{%
To frame the scope. The project proves a bracket around the robust radius, shows a dissociation
across model families, debunks the weak-attack effect, and adds a shape term for ranking robust
models. It does not claim the ratio is a new object, it does not claim co-existence of invariance and
robustness as a discovery, and it does not propose a new attack or a new defense. The ratio is a
diagnostic.}

\slide{5}{Where this started (practical-course recap)}{%
This began as the practical course, which I did with you. Two points. First, what we did: we
reproduced Ge's ordering with FGSM and PGD, then extended it to the standard architecture knobs,
pooling, kernel size, depth, a dense layer. Second, what we found and what it left open: the simple
``more consistency, less robust'' ordering broke, several CNNs were both consistent and as robust as
the fully-connected baseline, and we screened models on a consistency-versus-robustness graph. But it
never said which quantity predicts robustness when consistency does not, and it used only weak
attacks. That is what the rest of the project addresses.}

\slide{6}{From the project to now}{%
Two things happened since. First, the co-existence we saw in the course is now in a published paper,
TIPS from Saha and Gokhale in 2024, so it is established, not something we are claiming. Second, there
are still additions we can make on top of it: nobody ran a strong parameter-free attack on this line,
nobody checked whether the invariance metric actually orders robustness, and there was no attack-free
predictor of robustness. So the question is: is there such a quantity? The answer is yes, and it is a
margin quantity.}

\slide{7}{Invariance and robustness are different directions}{%
Geometrically these are different directions. On the left, imposing invariance projects the classes
onto their average part, which can shrink the margin; that is Ge's mechanism. Robustness is the
distance to the decision boundary. On the right, take one point: a shift moves it around the boundary,
so the label is kept; an attack moves it toward the boundary and crosses. The two directions are
nearly orthogonal, so a good invariance score need not buy you robustness.}

\slide{8}{The quantity, from a standard certificate}{%
The quantity itself is not new, it is the Lipschitz--margin certificate. The margin, divided by the
size of its gradient, lower-bounds the robust radius. We do two things to it. We match the norm to the
attack, the dual norm, which matters: the wrong norm flips the sign of the correlation. And we note it
is scale-free, so robustness attaches to the ratio, not to the margin or the Lipschitz constant
separately. We call the model-level summary eta over L, the mean margin over the mean dual-norm
gradient on clean correctly-classified images. No attack is needed to compute it.}

\slide{9}{Proposition 1: invariance fixes the usable margin}{%
Proposition one. A shift-invariant linear model must have a constant weight vector, so it depends only
on the average of the input. The best margin it can achieve is then half the distance between the two
class averages. In plain words: invariance throws away everything but the average, and the average is
close for the two classes, so the margin is small. Ge proved this for the single-image case; we prove
it for any finite separable dataset.}

\slide{10}{Proposition 2: the ratio brackets the true radius}{%
Proposition two. The certificate gives a lower bound on the radius. We prove a matching upper bound,
so the true radius is trapped between eta over L and eta over alpha. The gap between the two is a
condition number kappa, and we measure it between 1.6 and 2.4 in our tests, so the lower bound is
essentially tight and eta over L is essentially the radius. The lower bound is Hein and Tsuzuku; the
upper bound is what we add.}

\slide{11}{Does shift-consistency order robustness? Does the ratio?}{%
Here is the main dissociation. We build four invariance operators, standard, anti-aliased, exactly
cyclic, and shift-augmented, with identical parameter counts, so only the invariance mechanism
differs. Across three datasets and up to a ResNet-18, shift-consistency does not order robustness, and
under adversarial training it actually anti-predicts it. The threat-matched ratio orders it, and it
carries information beyond clean accuracy, the partial correlation is plus 0.89.}

\slide{12}{Is the ``invariance helps'' effect real?}{%
Is the ``helps'' effect real? We rerun that exact regime, weak attacks on undefended models, on our
operators and on Saha's and Wang's, and then we add AutoAttack. Under FGSM the ordering shows up.
Under AutoAttack every non-robust model drops to near zero and the ordering is gone. On frozen CLIP
encoders, FGSM reports 16 to 61 percent robust accuracy that AutoAttack erases. So the ``helps''
effect is a weak-attack artifact, and nobody in that line had run a strong attack.}

\slide{13}{Per image, where training cannot confound}{%
The cleanest test. On sixteen frozen CLIP encoders the invariance metric is saturated, it sits between
0.963 and 0.988 for everyone, robust or not, so it cannot rank them. So we go per image inside a
single encoder, where training cannot be the confound: there the ratio orders each image's radius
strongly, Spearman 0.74 to 0.96, and per-image consistency does not, it runs from 0.37 down to about
zero.}

\slide{14}{Ranking already-robust models: the gradient's shape}{%
Among already-robust models the ratio no longer separates them. What does is the shape of the
gradient, an L1-over-L2 anisotropy between 1 and root d. An infinity-norm attack wastes budget on a
concentrated gradient and exploits a spread-out one, so lower anisotropy is more robust. On thirty
RobustBench CIFAR-10 models the correlation is minus 0.79, and minus 0.77 controlling for clean
accuracy, and it replicates on MNIST and Fashion. We propose a spread-over-dimensions deficit as the
mechanism and verify it on a controlled model and thirty real ones, but we call it a proposed
mechanism, not a theorem.}

\slide{15}{The results across settings}{%
One table pulling it together. The threat-matched ratio orders robustness in every setting.
Shift-consistency does not, and near zero means it does not order the radius at all. The shape term
ranks the already-robust models, so it only appears in the rows where we rank robust models. The blank
cells are quantities we did not measure in that setting, for instance we did not compute
shift-consistency for the downloaded RobustBench models. Every number here reproduces from the code
and result files.}

\slide{16}{Where it breaks}{%
Where it breaks, stated on purpose. First, the ratio cannot be trained toward, it is scale-free and
its direct maximizer is the constant classifier, so it is a diagnostic, not a loss. Second, the shape
law needs a wide robustness range to show up; on low-dimensional data it is compressed. Third, the
deficit is a proposed mechanism, not a network-level proof. Fourth, on text the worst case over a
small paraphrase set is too mild to be a real second adversarial axis.}

\slide{17}{Proposed contributions, and existing results}{%
To be explicit about what is ours. On the left, the proposed contributions: the dissociation, the
powered per-image version inside one encoder, the weak-attack result, the upper arm of the bracket,
the anisotropy ranking and its mechanism, the saturation finding, and the sign-flip under adversarial
training. On the right, the existing results we use or refine: the certificate itself, CLEVER, the
fact these are related at all, Ge's DC mechanism, randomized smoothing, the dimension and curvature
views, the participation ratio, AutoAttack, and Ngnawe's margin consistency.}

\slide{18}{Summary and directions}{%
To summarize. Shift-consistency, the metric this literature tracks, does not order adversarial
robustness, and the ``helps'' effect vanishes under a strong attack. A threat-matched, scale-free
margin-to-Lipschitz ratio, which we complete to a two-sided bracket, orders robustness from small
classifiers to frozen CLIP encoders. And to rank already-robust models you need the shape of the
gradient. Next directions: other norms beyond L-infinity, the rich-regime optimization question, a
real text-adversarial axis, and scaling the anisotropy study beyond RobustBench.}

\slide{19}{Where your input would help most}{%
This last slide is really the point of sending it to you. Four things. First, positioning: eta over L
is a certificate and we add a two-sided bracket; given your certificate work, is that a genuine
addition or is it subsumed by smoothing and Lipschitz certification, and how would you frame
certificate versus diagnostic. Second, the kernel theory: our lazy-regime result treats orbit-averaging
as a projection in the tangent-kernel space; given your deep-net-as-kernel-machine work, is that the
right lens and where are the pitfalls. Third, which parts look weakest to you and what would you rerun
or design differently. Fourth, are we missing work at the certified-robustness-times-invariance
intersection, and which result should headline, the dissociation or the anisotropy law.}

\vspace{1.2em}
\hrule
\vspace{0.4em}
\section{Backup: anticipated questions}

{\small\color{slate} Not in the shared deck. Short answers to have ready.}

\subsection*{Novelty and positioning}
\qa{1. Is \etaL{} just the Lipschitz--margin bound?}{The lower bound, yes. We add threat-matching, the
scale-free reading, the upper bound, and the empirical dissociation across families.}
\qa{2. Saha/TIPS reports invariance helps. How do you differ?}{That link is a weak-attack reading; under
AutoAttack consistency does not order robustness.}
\qa{3. Ge already showed invariance can hurt.}{Ge gives a margin-collapse on specific data; we give the
ratio that predicts robustness across models, with Ge's condition as a corollary.}
\qa{4. Wang 2025 also uses the word ``anisotropy''.}{Theirs is on-orbit gradient variability for
equivariance; ours is a scalar $\ell_1/\ell_2$ ratio that ranks robust models. Different object; their
diagnostic tool is CLEVER, the same margin/Lipschitz family as ours.}
\qa{5. Is the anisotropy result just CLEVER or the participation ratio?}{CLEVER estimates a radius; the
participation ratio is for catastrophic-overfitting detection during training. Ranking robust models by
gradient spread is the new part.}
\qa{6. Why not just report robust accuracy?}{\etaL{} needs no attack, is per-sample, and ranks cheaply;
robust accuracy is the ground truth we validate against.}

\subsection*{Methodology}
\qa{7. Which attack is the arbiter?}{AutoAttack; we check gradient masking (AutoAttack $\le$ matched PGD)
wherever we claim robustness.}
\qa{8. Why does FGSM vs AutoAttack matter?}{FGSM over-reports robustness by a large factor; the ``helps''
ordering only survives FGSM.}
\qa{9. How is the radius measured?}{A per-sample minimum-norm attack (DDN), validated against the analytic
linear radius; not grid-PGD.}
\qa{10. Capacity-matched, really?}{The four arms have identical parameter counts; only the invariance
operator differs.}
\qa{11. Seeds and confidence intervals?}{Multiple seeds; bootstrap CIs and permutation tests on the dense
grids.}
\qa{12. Is BatchNorm a confound (Galloway)?}{Controlled for; the dissociation survives BN-free and
matched-capacity arms.}
\qa{13. Threat-matching, concretely?}{$\ell_\infty$ uses the $\ell_1$ gradient norm, $\ell_2$ the $\ell_2$
norm; the mismatched norm gives the wrong sign on MNIST.}
\qa{14. Clean-accuracy confound?}{Partial correlations control for it (\etaL{} vs AutoAttack partial
$+0.89$; anisotropy partial $-0.77$).}

\subsection*{Theory and scope}
\qa{15. Is the bracket tight?}{Exact ($\kappa=1$) for linear invariant features; $\kappa\in[1.6,2.4]$ for
the power-spectrum surrogate.}
\qa{16. What is $\alpha$?}{The slowest rate the margin can fall along a path to the boundary, a
co-Lipschitz slope.}
\qa{17. Does Proposition 1 need orbit-closure?}{No; it holds for any finite separable dataset via
projection onto the invariant subspace.}
\qa{18. Is the deficit a theorem?}{No; a proposed mechanism, verified on a controlled model and 30 real
ones.}
\qa{19. Why is curvature not the ranker?}{On-direction curvature does not rank robust models and the
anisotropy does not track it; the deficit is transverse.}
\qa{20. Relation to Simon-Gabriel?}{They fix equal per-coordinate variance, so $A\propto\sqrt d$; we let
$A$ vary at fixed $d$ and it tracks robustness.}
\qa{21. Can \etaL{} be trained?}{No; scale-free, so its direct maximizer is the constant classifier.}
\qa{22. What is the practical use?}{A cheap, attack-free screen for architectures before paying for
adversarial training.}
\qa{23. What would falsify the claim?}{A family where consistency orders threat-matched robustness and
\etaL{} does not; we did not find one.}

\end{document}
